\documentclass[11pt,a4paper]{article}
\usepackage[utf8]{inputenc}
\usepackage[T1]{fontenc}
\usepackage[english]{babel}
\usepackage{geometry}
\usepackage{amssymb}
\usepackage{parskip}
\usepackage{mathptmx}

\geometry{margin=2.5cm}

\begin{document}

\begin{titlepage}
\thispagestyle{empty}
\centering
\vspace*{2cm}
T.C.\\
IZMIR UNIVERSITY OF ECONOMICS\\
FACULTY OF ENGINEERING\\
CE 216
\vfill
{\Large\bfseries ``SPORTS MANAGER'' PROJECT}\\[0.75em]
{\large\bfseries MILESTONE 3 REPORT DOCUMENT}
\vfill
SİMGE GÖKALP — 20230602029\\[0.45em]
TUANA YAĞMUR — 20240602410\\[0.45em]
VİLDAN TURNALAR İNAL — 20253501001\\[0.45em]
UĞUR EMİN BAYNAL — 20220602408\\[1.25em]
LECTURER: KAYA OĞUZ
\end{titlepage}

\section{Introduction}

This report presents the final implementation of the Sports Manager application developed during the CE216 course project. The project began with an architecture-focused design document in Milestone 1, continued with the implementation of the core framework and testing infrastructure in Milestone 2, and was completed in Milestone 3 with a fully playable application, graphical user interface, second sport integration, installer-ready structure, and user-oriented features.

The primary objective of the project was to create a sports management framework capable of supporting multiple sports through abstraction and object-oriented design principles. Instead of creating a football-only game, the system was designed as a reusable and extensible framework where different sports can share the same core architecture while implementing their own rules and mechanics.

The final system successfully demonstrates this goal. The application now supports two sports:
\begin{itemize}
\item Football
\item Basketball
\end{itemize}
Both sports operate on the same framework while implementing different match logic, player structures, tactics, lineups, and league behavior.

The project includes:
\begin{itemize}
\item A JavaFX graphical user interface
\item Match simulation systems
\item League and standings management
\item Squad and tactic management
\item Injury and substitution systems
\item Theme and settings support
\item Save-ready architecture
\item Multi-sport extensibility
\item Unit testing with JUnit
\end{itemize}

The final implementation preserves the layered architecture and abstraction principles defined during Milestone 1 while introducing several practical refinements based on implementation experience.

\section{Final System Architecture}

The Sports Manager application follows a layered and modular architecture designed to separate responsibilities and minimize coupling between components.

\subsection{User Interface Layer}

The User Interface layer is implemented using JavaFX and FXML.

Controllers include:
\begin{itemize}
\item MainMenuController
\item DashboardController
\item FixtureController
\item MatchScreenController
\item LeagueTableController
\item SquadController
\item TeamSelectionController
\item SportSelectionController
\item SettingsController
\item AppToolbarController
\end{itemize}

The UI layer never communicates directly with concrete sport implementations. Controllers interact only with abstract types and the \texttt{SportManager} class.

Scene transitions are centralized inside the \texttt{SceneManager} class. This prevents duplication of navigation logic and keeps controllers lightweight.

\subsection{Core System Layer}

The Core System layer contains the abstract framework shared by all sports.

Main abstract classes:
\begin{itemize}
\item Player
\item Team
\item Match
\item League
\end{itemize}

Supporting classes:
\begin{itemize}
\item StandingEntry
\item Lineup
\item Tactic
\item MatchResult
\item MatchSegment
\item InjuryRecord
\item Coach
\end{itemize}

The \texttt{SportManager} class acts as the facade between the UI and the game system.

\texttt{GameSession} is responsible for maintaining application state including:
\begin{itemize}
\item Current season year
\item Current round/week
\item Managed team
\item Selected sport
\item Active league
\item Match progression
\end{itemize}

\subsection{Sport Abstraction Layer}

The \texttt{Sport} interface defines all sport-specific behavior.

Implemented methods include:
\begin{itemize}
\item \texttt{getName()}
\item \texttt{createLeague()}
\item \texttt{createTeam()}
\item \texttt{createMatch()}
\item \texttt{getPositions()}
\item \texttt{getTactics()}
\item \texttt{getRequiredLineupSize()}
\item \texttt{getMaxSubstituteCount()}
\item \texttt{getSegmentCount()}
\item \texttt{getSegmentLabel()}
\end{itemize}

\texttt{SportFactory} centralizes concrete sport creation logic.

Example:
\begin{verbatim}
case "football":
    return new FootballSport();
case "basketball":
    return new BasketballSport();
\end{verbatim}

This ensures that the rest of the application remains independent from sport implementations.

\subsection{Concrete Sport Modules}

\subsubsection*{Football Module}

Implemented football classes:
\begin{itemize}
\item FootballSport
\item FootballLeague
\item FootballTeam
\item FootballPlayer
\item FootballMatch
\end{itemize}

Football uses:
\begin{itemize}
\item 11-player lineups
\item 2 halves
\item Football formations
\item Goal-based scoring
\item Standard 3-1-0 point system
\end{itemize}

\subsubsection*{Basketball Module}

Implemented basketball classes:
\begin{itemize}
\item BasketballSport
\item BasketballLeague
\item BasketballTeam
\item BasketballPlayer
\item BasketballMatch
\item BasketballTactics
\end{itemize}

Basketball introduces:
\begin{itemize}
\item Quarter-based matches
\item Basketball-specific tactics
\item Higher-scoring simulations
\item Smaller lineups
\item Faster possession-based gameplay
\end{itemize}

The successful integration of basketball demonstrates that the framework supports extension without requiring major architectural changes.

\section{Object-Oriented Design Principles Applied}

The project heavily relies on object-oriented programming principles.

\subsection{Abstraction}

Abstraction is the foundation of the entire project.

Examples:
\begin{itemize}
\item Abstract classes: Player, Team, League, Match
\item Sport interface
\item Generic UI interaction through abstract types
\end{itemize}

The UI does not know whether it is managing football or basketball. It communicates only through abstract interfaces.

\subsection{Encapsulation}

Data and behavior are encapsulated inside domain objects.

Examples:
\begin{itemize}
\item Player manages injury recovery internally
\item Team manages lineup validation
\item League controls standings and rounds
\item Match controls simulation flow
\end{itemize}

Sensitive fields are protected using private/protected access modifiers and accessed through getters/setters.

\subsection{Inheritance}

Concrete sports inherit shared functionality from abstract base classes.

Examples:
\begin{itemize}
\item FootballPlayer extends Player
\item BasketballPlayer extends Player
\item FootballTeam extends Team
\item BasketballLeague extends League
\end{itemize}

This reduces code duplication while allowing sport-specific behavior.

\subsection{Polymorphism}

Polymorphism allows the application to work with multiple sports through shared interfaces.

Example:
\begin{verbatim}
Sport sport = SportFactory.create(sportCode);
\end{verbatim}

The caller does not need to know whether the returned object is football or basketball.

\subsection{Factory Pattern}

\texttt{SportFactory} centralizes object creation.

Benefits:
\begin{itemize}
\item Reduced coupling
\item Easy extension
\item Cleaner UI layer
\item Simplified initialization
\end{itemize}

\subsection{Separation of Concerns}

Responsibilities are clearly separated:
\begin{itemize}
\item UI handles presentation
\item SportManager coordinates actions
\item GameSession stores state
\item Domain objects manage game logic
\item Sport modules implement specific rules
\end{itemize}

\section{Changes From Milestone 1 Design}

Several improvements and practical refinements were introduced during implementation.

\subsection{Match System Redesign}

The original start/play/finish structure was replaced by a segment-based simulation system.

Benefits:
\begin{itemize}
\item Better UI integration
\item Easier halftime management
\item Better event generation
\item Cleaner simulation flow
\end{itemize}

\subsection{Team Statistics Refactoring}

Season statistics were moved directly into Team instead of being duplicated inside StandingEntry.

Benefits:
\begin{itemize}
\item Eliminated synchronization issues
\item Reduced duplicated data
\item Simplified standings updates
\end{itemize}

\subsection{Simplified Tactic System}

The original design included numeric attack/defense levels.

The final implementation simplified tactics into:
\begin{itemize}
\item Name
\item Style
\end{itemize}

This was sufficient for the implemented simulation systems while keeping the design extensible.

\subsection{SceneManager Addition}

\texttt{SceneManager} was added to centralize JavaFX navigation.

Benefits:
\begin{itemize}
\item Cleaner controllers
\item Better separation of concerns
\item Easier screen transitions
\end{itemize}

\subsection{Basketball Integration}

The biggest architectural extension from Milestone 2 to Milestone 3 was the addition of basketball.

This required:
\begin{itemize}
\item New sport implementation
\item New player types
\item Basketball-specific match simulation
\item Possession-based logic
\item Quarter system
\item Basketball tactics
\item Basketball standings
\end{itemize}

The integration validated the flexibility of the framework.

\section{Implemented Features}

\subsection{League and Season System}

The application generates fully playable league seasons.

Features:
\begin{itemize}
\item Automatic fixture generation
\item Double round-robin scheduling
\item Round progression
\item League standings
\item Champion determination
\item Week management
\end{itemize}

\subsection{Match Simulation}

\subsubsection*{Football}

Football simulation includes:
\begin{itemize}
\item Two halves
\item Tactical modifiers
\item Goals
\item Yellow cards
\item Injuries
\item Substitutions
\item Live event feed
\end{itemize}

\subsubsection*{Basketball}

Basketball simulation includes:
\begin{itemize}
\item Quarter-based gameplay
\item Possession simulation
\item Basketball tactics
\item Offensive and defensive balancing
\item Fast scoring system
\end{itemize}

\subsection{Tactical System}

The tactical system supports multiple formations and styles.

Football formations:
\begin{itemize}
\item 4-3-3
\item 4-4-2
\item 4-2-3-1
\item 3-5-2
\item 5-3-2
\end{itemize}

Basketball tactics include offensive strategy configurations.

\subsection{Squad and Lineup Management}

Users can:
\begin{itemize}
\item View players
\item Create lineups
\item Change tactics
\item Make substitutions
\item Track injuries
\item View player attributes
\end{itemize}

\subsection{Injury System}

Players can become injured during matches.

Features:
\begin{itemize}
\item Random injury generation
\item Match-based recovery system
\item Injury frequency settings
\item Availability checks
\end{itemize}

\subsection{Settings System}

Implemented settings:
\begin{itemize}
\item Match difficulty
\item Injury frequency
\item Maximum substitutions
\item Accent themes
\item Auto-advance
\item Detailed event display
\end{itemize}

\subsection{User Interface}

The JavaFX interface includes:
\begin{itemize}
\item Animated splash screen
\item Dashboard
\item Fixtures screen
\item Match screen
\item Squad screen
\item League table
\item Settings page
\item Theme support
\item Toolbar navigation
\item Live match events
\end{itemize}

\section{Save and Load System}

The project architecture supports save/load functionality through Java serialization.

Core classes implement \texttt{Serializable}, including:
\begin{itemize}
\item Team
\item Player
\item League
\item Match
\item StandingEntry
\item Lineup
\item Tactic
\end{itemize}

This allows complete game state persistence.

The final implementation includes a functional save/load system based on serialized game snapshots, while still allowing future expansion for more advanced persistence features.

\section{Testing}

JUnit was used for testing.

Tested systems include:
\begin{itemize}
\item Team statistics
\item Standings sorting
\item Injury handling
\item Lineup validation
\item League progression
\item Sport factory creation
\item Match logic
\end{itemize}

The tests validate both:
\begin{itemize}
\item Framework behavior
\item Sport-specific implementations
\end{itemize}

The testing process significantly improved project reliability.

\section{Running the Project}

\subsection*{Requirements}

\begin{itemize}
\item Java JDK 21+
\item Apache Maven 3.8+
\end{itemize}

\subsection*{Build}

\begin{verbatim}
mvn compile
\end{verbatim}

\subsection*{Run}

\begin{verbatim}
mvn compile exec:java
\end{verbatim}

\subsection*{Run Tests}

\begin{verbatim}
mvn test
\end{verbatim}

\section{Post-Mortem}

\subsection{What Went Well}

Several aspects of the project were successful.

\subsubsection*{Strong Core Architecture}

The abstraction-focused design created during Milestone 1 proved highly effective during implementation.

The addition of basketball required only minimal changes outside the new module itself, demonstrating that the architecture was truly extensible.

\subsubsection*{Successful Multi-Sport Support}

The project successfully supports two sports using the same framework.

This validates:
\begin{itemize}
\item Interface-based design
\item Factory pattern usage
\item Shared domain abstraction
\item Polymorphic behavior
\end{itemize}

\subsubsection*{UI Development}

The JavaFX interface became significantly more polished during Milestone 3.

The addition of:
\begin{itemize}
\item Theme support
\item Toolbar navigation
\item Animated transitions
\item Match event feeds
\item Tactical visuals
\end{itemize}

improved overall usability and presentation quality.

\subsubsection*{Team Collaboration}

Git-based collaboration and task sharing allowed the team to work on different modules simultaneously.

Responsibilities were divided across:
\begin{itemize}
\item UI
\item Match simulation
\item Core framework
\item Testing
\item Basketball integration
\end{itemize}

\subsection{Challenges and Problems}

\subsubsection*{JavaFX and Maven Configuration}

One of the biggest technical difficulties involved JavaFX runtime configuration on Windows systems.

Problems included:
\begin{itemize}
\item Module-path issues
\item Native library loading
\item Non-ASCII path problems
\item Maven plugin configuration
\end{itemize}

Several workarounds were required to ensure stable execution.

\subsubsection*{Merge Conflicts}

As the project grew larger, merge conflicts became more common.

This was especially noticeable when:
\begin{itemize}
\item Refactoring the tactic system
\item Modifying abstract classes
\item Updating Team and Match logic
\end{itemize}

The team improved its Git workflow over time to reduce conflicts.

\subsubsection*{Scope Management}

Some planned features were simplified or partially implemented due to time constraints.

Examples:
\begin{itemize}
\item Advanced transfer system
\item Financial management
\item Detailed training mechanics
\item Full save/load UI
\item Multi-season persistence
\end{itemize}

\subsection{Lessons Learned}

The project provided practical experience in:
\begin{itemize}
\item Object-oriented architecture
\item JavaFX development
\item Maven dependency management
\item Git collaboration
\item Large-scale Java project structure
\item Abstraction and extensibility
\item Refactoring
\item Team coordination
\end{itemize}

The importance of designing extensible systems early became especially clear during basketball integration.

\section{Final Evaluation}

The Sports Manager project successfully evolved from a conceptual design into a fully functional multi-sport management application.

The final system demonstrates:
\begin{itemize}
\item Strong object-oriented design
\item Functional abstraction
\item Multi-sport support
\item Modular architecture
\item Working graphical interface
\item Extendable framework
\item Stable gameplay systems
\end{itemize}

The project successfully satisfies the core objectives defined in the assignment specification.

The implementation proves that the framework can continue growing with additional sports and features in future iterations.

\end{document}
